\SourcePage{71}{74}
\section{The Schrödinger equation}

The Hamiltonian determines the form of a physical system's wave equation and
consequently acquires fundamental significance in the entire mathematical
apparatus of quantum mechanics.

For a free particle, the Hamiltonian is already established by the general
requirements associated with the homogeneity and isotropy of space and with
the Galilean relativity principle. In classical mechanics these requirements
imply a quadratic dependence of the particle energy on its
momentum: $E=p^2/2m$, where the constant $m$ is called the mass of the particle
(see I, \S~4). In quantum mechanics the same requirements lead to the same
relation for the eigenvalues of energy and momentum, quantities which for a
free particle are conserved and can be measured simultaneously.

For the relation $E=p^2/2m$ to hold for every eigenvalue of energy and
momentum, however, it must also hold for their operators:
\begin{equation}
  \hat H=\frac{1}{2m}\left(\hat p_x^2+\hat p_y^2+\hat p_z^2\right).
  \tag{17.1}
\end{equation}
Substitution of (15.2) here gives the Hamiltonian of a freely moving particle
in the form
\begin{equation}
  \hat H=-\frac{\hbar^2}{2m}\Delta,
  \tag{17.2}
\end{equation}
where $\Delta=\partial^2/\partial x^2+\partial^2/\partial y^2+
\partial^2/\partial z^2$ is the Laplace operator. The Hamiltonian of a system
of noninteracting particles is equal to the sum of the Hamiltonians of the
individual particles:
\begin{equation}
  \hat H=-\frac{\hbar^2}{2}\sum_a\frac{\Delta_a}{m_a},
  \tag{17.3}
\end{equation}
where the index $a$ numbers the particles; $\Delta_a$ is the Laplace operator
in which differentiation is performed with respect to the coordinates of the
$a$th particle.
\SourcePage{72}{75}
In classical (nonrelativistic) mechanics, particle interaction is described by
an additive term in the Hamilton function, the potential energy of interaction
$U(\mathbf r_1,\mathbf r_2,\ldots)$, which is a function of the particle
coordinates. Particle interaction in quantum mechanics is likewise described
by adding the same function to the Hamiltonian of the
system\footnote{This assertion is not, of course, a logical consequence of the
fundamental principles of quantum mechanics and must be regarded as a
consequence of experimental data.}:
\begin{equation}
  \hat H=-\frac{\hbar^2}{2}\sum_a\frac{\Delta_a}{m_a}
  +U(\mathbf r_1,\mathbf r_2,\ldots);
  \tag{17.4}
\end{equation}
the first term may be regarded as the kinetic-energy operator and the second
as the potential-energy operator. In particular, the Hamiltonian for one
particle in an external field is
\begin{equation}
  \hat H=\frac{\hat{\mathbf p}^{2}}{2m}+U(x,y,z)
  =-\frac{\hbar^2}{2m}\Delta+U(x,y,z),
  \tag{17.5}
\end{equation}
where $U(x,y,z)$ is the potential energy of the particle in the external
field.

Substitution of expressions (17.2)--(17.5) in the general equation (8.1)
yields the wave equations for the corresponding systems. Here we write out
the wave equation for a particle in an external field:
\begin{equation}
  i\hbar\frac{\partial\Psi}{\partial t}
  =-\frac{\hbar^2}{2m}\Delta\Psi+U(x,y,z)\Psi.
  \tag{17.6}
\end{equation}
Equation (10.2), determining the stationary states, assumes the form
\begin{equation}
  \frac{\hbar^2}{2m}\Delta\psi+[E-U(x,y,z)]\psi=0.
  \tag{17.7}
\end{equation}

Equations (17.6) and (17.7) were established by \emph{Schrödinger} in 1926 and
are known as the \emph{Schrödinger equations}.

For a free particle equation (17.7) becomes
\begin{equation}
  \frac{\hbar^2}{2m}\Delta\psi+E\psi=0.
  \tag{17.8}
\end{equation}
This equation admits solutions that remain finite over all space for any positive value
of the energy $E$. In states with definite directions of motion these
solutions are the eigenfunctions of the momentum operator, with%
\SourcePage{73}{76}
\relax\ $E=p^2/2m$. The full (time-dependent) wave functions of these
stationary states have the form
\begin{equation}
  \Psi=\operatorname{const}\cdot
  \exp\left[-\frac{i}{\hbar}(Et-\mathbf p\cdot\mathbf r)\right].
  \tag{17.9}
\end{equation}
Each of these functions, a \emph{plane wave}, describes a state in which the
particle possesses definite energy $E$ and momentum $\mathbf p$. The frequency of
this wave is $E/\hbar$, and its wave vector is
$\mathbf k=\mathbf p/\hbar$; the associated wavelength
$\lambda=2\pi\hbar/p$ is termed the \emph{de Broglie wavelength of the
particle}\footnote{The concept of a wave associated with a particle was first
introduced by \emph{de Broglie} (\emph{L.~de Broglie}) in 1924.}.

The energy spectrum of a freely moving particle is therefore continuous,
extending from zero to $+\infty$. Each of these eigenvalues, except only
$E=0$, is degenerate, and the degeneracy is of infinite multiplicity. Indeed,
every nonzero value of $E$ corresponds to an infinite set of eigenfunctions
(17.9), differing in the directions of the vector $\mathbf p$ while its
absolute magnitude remains the same.

Let us trace how the limiting transition to classical mechanics occurs in the
Schrödinger equation, considering for simplicity just one particle in an
external field. Substitution in the Schrödinger equation (17.6) of the limiting
expression (6.1), $\Psi=ae^{iS/\hbar}$, for the wave function gives, after the
differentiations are performed,
\[
  a\frac{\partial S}{\partial t}
  -i\hbar\frac{\partial a}{\partial t}
  +\frac{a}{2m}{\left(\nabla S\right)}^2
  -\frac{i\hbar}{2m}a\Delta S
  -\frac{i\hbar}{m}\nabla S\cdot\nabla a
  -\frac{\hbar^2}{2m}\Delta a+Ua=0.
\]
This equation contains purely real and purely imaginary terms (recall that $S$
and $a$ are real); equating each group separately to zero gives two equations:
\[
  \frac{\partial S}{\partial t}
  +\frac{1}{2m}{\left(\nabla S\right)}^2
  +U-\frac{\hbar^2}{2ma}\Delta a=0,
\]
\[
  \frac{\partial a}{\partial t}
  +\frac{a}{2m}\Delta S+\frac{1}{m}\nabla S\cdot\nabla a=0.
\]
Neglecting the term containing $\hbar^2$ in the first of these equations, we
obtain
\begin{equation}
  \frac{\partial S}{\partial t}
  +\frac{1}{2m}{\left(\nabla S\right)}^2+U=0,
  \tag{17.10}
\end{equation}
that is, as expected, the classical Hamilton--Jacobi equation for the action
$S$ of the particle. We note in passing that, as $\hbar\to0$%
\SourcePage{74}{77}
\relax\ classical mechanics is accurate up to and including quantities of
first order in $\hbar$ (rather than zeroth order).

After multiplication by $2a$, the second equation obtained above can be
rewritten as
\begin{equation}
  \frac{\partial a^2}{\partial t}
  +\Div{\left(a^2\frac{\nabla S}{m}\right)}=0.
  \tag{17.11}
\end{equation}
This equation has a transparent physical meaning: $a^2$ is the probability
density for finding the particle at one or another place in space
($|\Psi|^2=a^2$); $\nabla S/m=\mathbf p/m$ is the classical velocity
$\mathbf v$ of the particle. Equation (17.11) is therefore nothing other than
the continuity equation, showing that the probability density is
``transported'' according to the laws of classical mechanics, with the
classical velocity $\mathbf v$ at every point.

\ProblemHeading
Find the transformation law of the wave function under a Galilean
transformation.

\SolutionHeading
Perform the transformation on the wave function for the free motion of a
particle (a plane wave). Since every function $\Psi$ can be expanded in plane
waves, this will also determine the transformation law for an arbitrary wave
function.

The plane waves in the reference frames $K$ and $K'$ (where $K'$ moves with
velocity $\mathbf V$ relative to $K$) are
\[
  \Psi(\mathbf r,t)=\operatorname{const}\cdot
  \exp[i(\mathbf p\cdot\mathbf r-Et)/\hbar],
\]
\[
  \Psi'(\mathbf r',t)=\operatorname{const}\cdot
  \exp[i(\mathbf p'\cdot\mathbf r'-E't)/\hbar],
\]
where $\mathbf r=\mathbf r'+\mathbf Vt$, while the particle's momenta and
energies in the two frames are related by
\[
  \mathbf p=\mathbf p'+m\mathbf V,
  \qquad
  E=E'+\mathbf V\cdot\mathbf p'+\frac{mV^2}{2}
\]
(see I, \S~8). Substitution of these expressions in $\Psi$ gives
\begin{equation}
  \begin{aligned}
    \Psi(\mathbf r,t)
      &=\Psi'(\mathbf r',t)
        \exp\left[\frac{i}{\hbar}
          \left(m\mathbf V\cdot\mathbf r'+\frac{mV^2}{2}t\right)\right]\\
      &=\Psi'(\mathbf r-\mathbf Vt,t)
        \exp\left[\frac{i}{\hbar}
          \left(m\mathbf V\cdot\mathbf r-\frac{mV^2}{2}t\right)\right].
  \end{aligned}
  \tag{1}
\end{equation}
Written this way, the formula no longer contains quantities characterizing
the free motion of the particle; it therefore furnishes the sought-after
general rule governing how the wave function of an arbitrary state of the
particle transforms. For a system of particles, the exponent of the
exponential in (1) must contain a summation over all the particles.
